\documentclass[conference]{IEEEtran}
\usepackage{amsmath,amssymb,amsfonts}
\usepackage{algorithmic}
\usepackage{graphicx}
\usepackage{textcomp}
\usepackage{xcolor}
\usepackage{booktabs}
\usepackage{cite}
\usepackage{microtype}
\usepackage{tabularx}
\usepackage{url}
\usepackage{balance}

\begin{document}

\title{Generative Adversarial Networks: Mathematical Foundations, Minimax Game Theory, Architectural Evolution, and Empirical Frontiers}

\maketitle

\begin{abstract}
Generative Adversarial Networks (GANs) represent a foundational breakthrough in unsupervised and self-supervised deep learning, formulating generative modeling as a minimax zero-sum game between two competing neural networks: a Generator and a Discriminator. This paper presents an exhaustive academic investigation into GANs, spanning rigorous mathematical derivations from first principles, training dynamics, failure modes, architectural evolutions, and real-world applications. We derive the optimal discriminator formulation, prove that the global optimum of the minimax objective minimizes the Jensen-Shannon divergence between the real and generated probability measures, and analyze the conditions under which a Nash equilibrium is attainable. Furthermore, we conduct an in-depth comparative analysis across four seminal milestone papers: the original GAN formulation (Goodfellow et al., 2014), Deep Convolutional GANs (Radford et al., 2015), Wasserstein GANs with Gradient Penalty (Arjovsky et al., 2017; Gulrajani et al., 2017), and Style-based GANs (Karras et al., 2019/2020). Finally, we categorize the practical deployment of GANs across computer vision, super-resolution, cross-domain image synthesis, and medical imaging, concluding with a formal evaluation of quantitative fidelity metrics including Inception Score (IS) and Fr\'echet Inception Distance (FID).
\end{abstract}

\begin{IEEEkeywords}
Generative Adversarial Networks, Minimax Optimization, Jensen-Shannon Divergence, Wasserstein Distance, DCGAN, WGAN-GP, StyleGAN, Fr\'echet Inception Distance.
\end{IEEEkeywords}

\section{Introduction}
Generative modeling in machine learning aims to approximate an underlying data distribution $p_{data}$ given observed samples $x \in \mathcal{X}$, such that new plausible synthetic samples can be generated. Prior to 2014, primary approaches fell into explicit density estimation frameworks, such as tractable autoregressive models or intractable approximations requiring Markov Chain Monte Carlo (MCMC) sampling or variational lower bounds (e.g., Variational Autoencoders).

In 2014, Goodfellow et al. \cite{goodfellow2014generative} introduced Generative Adversarial Networks (GANs), fundamentally bypassing explicit density estimation. Instead, GANs re-conceptualize generative modeling as a game-theoretic contest between two parameterized differentiable functions:
\begin{itemize}
    \item \textbf{The Generator ($G$):} A mapping $G(z; \theta_g): \mathcal{Z} \to \mathcal{X}$ from a low-dimensional stochastic latent prior $p_z(z)$ into data space $\mathcal{X}$, acting as a counterfeiter striving to synthesize authentic instances.
    \item \textbf{The Discriminator ($D$):} A mapping $D(x; \theta_d): \mathcal{X} \to [0, 1]$, outputting the scalar probability that an input $x$ originates from true empirical data $p_{data}$ rather than the synthetic distribution $p_g$.
\end{itemize}

Through alternating gradient descent, the two networks engage in a two-player zero-sum game. When optimal, the generator recovers the true data distribution $p_g = p_{data}$, rendering the discriminator unable to distinguish between real and synthesized samples.

\section{Mathematical Foundations and Derivations}

\subsection{The Minimax Objective Function}
Let $x \sim p_{data}(x)$ denote real samples defined on $\mathbb{R}^d$, and let $z \sim p_z(z)$ denote latent noise vectors drawn from a known prior distribution (typically a standard isotropic Gaussian $\mathcal{N}(0, I)$ or uniform distribution $\mathcal{U}[-1, 1]$). The adversarial interaction is formalized as the two-player minimax game:
\begin{equation}
\min_G \max_D V(D, G)
\end{equation}
where the value function $V(D, G)$ is defined as:
\begin{align}
V(D, G) = &\; \mathbb{E}_{x \sim p_{data}} [\log D(x)] \nonumber \\
&+ \mathbb{E}_{z \sim p_z} [\log (1 - D(G(z)))]
\label{eq:minimax}
\end{align}
Expressed continuously over integrals:
\begin{align}
V(D, G) = &\int_{\mathcal{X}} p_{data}(x) \log D(x) \, dx \nonumber \\
&+ \int_{\mathcal{Z}} p_z(z) \log(1 - D(G(z))) \, dz
\end{align}
By applying the change of variables using the push-forward measure $p_g(x)$ induced by $G(z)$:
\begin{equation}
\int_{\mathcal{Z}} p_z(z) \log(1 - D(G(z))) \, dz = \int_{\mathcal{X}} p_g(x) \log(1 - D(x)) \, dx
\end{equation}
Thus, the continuous objective can be unified under a single integral over $\mathcal{X}$:
\begin{equation}
V(D, G) = \int_{\mathcal{X}} \left[ p_{data}(x) \log D(x) + p_g(x) \log(1 - D(x)) \right] dx
\label{eq:integral_unified}
\end{equation}

\subsection{Derivation of the Optimal Discriminator $D^*(x)$}
To determine the optimal discriminator $D^*$ for any fixed generator $G$, we maximize the integrand of (\ref{eq:integral_unified}) pointwise. Let $y = D(x) \in (0, 1)$, $a = p_{data}(x) \ge 0$, and $b = p_g(x) \ge 0$. We define the objective function:
\begin{equation}
f(y) = a \log y + b \log(1 - y)
\end{equation}
Taking the first derivative with respect to $y$:
\begin{equation}
\frac{df(y)}{dy} = \frac{a}{y} - \frac{b}{1 - y}
\end{equation}
Setting the first derivative to zero:
\begin{equation}
\frac{a}{y^*} - \frac{b}{1 - y^*} = 0 \implies a(1 - y^*) = b y^* \implies y^* = \frac{a}{a + b}
\end{equation}
Examining the second derivative:
\begin{equation}
\frac{d^2 f(y)}{dy^2} = -\frac{a}{y^2} - \frac{b}{(1 - y)^2} < 0, \quad \forall y \in (0, 1)
\end{equation}
Because the second derivative is strictly negative everywhere on $(0, 1)$, $f(y)$ is strictly concave, confirming that $y^*$ achieves the unique global maximum. Substituting back $a = p_{data}(x)$ and $b = p_g(x)$:
\begin{equation}
\boxed{D^*(x) = \frac{p_{data}(x)}{p_{data}(x) + p_g(x)}}
\label{eq:optimal_d}
\end{equation}

\subsection{Reformulation via Jensen-Shannon Divergence}
Having determined $D^*(x)$, we substitute (\ref{eq:optimal_d}) into the value function $V(D^*, G)$:
\begin{align}
V(D^*, G) = &\int_{\mathcal{X}} p_{data}(x) \log \left( \frac{p_{data}(x)}{p_{data}(x) + p_g(x)} \right) dx \nonumber \\
&+ \int_{\mathcal{X}} p_g(x) \log \left( \frac{p_g(x)}{p_{data}(x) + p_g(x)} \right) dx
\end{align}
Multiplying the denominators by $\frac{1}{2}$ inside the logarithms and subtracting $\log 2$ from both terms:
\begin{align}
V(D^*, G) = &\int_{\mathcal{X}} p_{data}(x) \left[ \log \left( \frac{p_{data}(x)}{\frac{p_{data}(x) + p_g(x)}{2}} \right) - \log 2 \right] dx \nonumber \\
&+ \int_{\mathcal{X}} p_g(x) \left[ \log \left( \frac{p_g(x)}{\frac{p_{data}(x) + p_g(x)}{2}} \right) - \log 2 \right] dx
\end{align}
Combining the constant terms:
\begin{align}
V(D^*, G) = &-\log 2 \int p_{data}(x) dx - \log 2 \int p_g(x) dx \nonumber \\
&+ D_{KL}\left( p_{data} \,\Big\|\, \frac{p_{data} + p_g}{2} \right) \nonumber \\
&+ D_{KL}\left( p_g \,\Big\|\, \frac{p_{data} + p_g}{2} \right)
\end{align}
Because probability density functions integrate to 1 ($\int p(x) dx = 1$):
\begin{align}
V(D^*, G) &= -2 \log 2 + D_{KL}\left( p_{data} \,\Big\|\, \frac{p_{data} + p_g}{2} \right) + D_{KL}\left( p_g \,\Big\|\, \frac{p_{data} + p_g}{2} \right) \nonumber \\
&= -\log 4 + 2 \cdot D_{JS}(p_{data} \parallel p_g)
\label{eq:v_jsd}
\end{align}
where $D_{JS}$ is the symmetric Jensen-Shannon divergence:
\begin{equation}
D_{JS}(P \parallel Q) = \frac{1}{2} D_{KL}\left(P \,\Big\|\, \frac{P+Q}{2}\right) + \frac{1}{2} D_{KL}\left(Q \,\Big\|\, \frac{P+Q}{2}\right)
\end{equation}

\subsection{Global Optimum and Nash Equilibrium}
The Jensen-Shannon divergence satisfies $D_{JS}(P \parallel Q) \ge 0$, with equality holding if and only if $P = Q$ almost everywhere. Consequently:
\begin{equation}
\min_G V(D^*, G) = -\log 4 \approx -1.3863
\end{equation}
which is uniquely achieved when:
\begin{equation}
p_g(x) = p_{data}(x), \quad \forall x \in \mathcal{X}
\end{equation}
At this global optimum, the discriminator output becomes completely uninformative:
\begin{equation}
D^*(x) = \frac{p_{data}(x)}{p_{data}(x) + p_{data}(x)} = \frac{1}{2}
\end{equation}
This state corresponds to the \textbf{Nash Equilibrium} of the zero-sum game in continuous function space: neither player can unilaterally decrease their loss by deviating from this strategy.

\section{Training Dynamics, Gradients, and Failure Modes}

\subsection{Alternating Stochastic Gradient Descent}
In practical optimization, $D$ and $G$ are parameterized by deep neural network weights $\theta_d$ and $\theta_g$. Optimization proceeds via alternating mini-batch updates:
\begin{algorithmic}[1]
\FOR{number of training iterations}
    \FOR{$k$ steps}
        \STATE Sample $\{z^{(1)}, \dots, z^{(m)}\} \sim p_z(z)$
        \STATE Sample $\{x^{(1)}, \dots, x^{(m)}\} \sim p_{data}(x)$
        \STATE Update $D$ by ascending stochastic gradient:
        $$\nabla_{\theta_d} \frac{1}{m} \sum_{i=1}^m \left[ \log D(x^{(i)}) + \log(1 - D(G(z^{(i)}))) \right]$$
    \ENDFOR
    \STATE Sample $\{z^{(1)}, \dots, z^{(m)}\} \sim p_z(z)$
    \STATE Update $G$ by descending stochastic gradient:
    $$\nabla_{\theta_g} \frac{1}{m} \sum_{i=1}^m \log(1 - D(G(z^{(i)})))$$
\ENDFOR
\end{algorithmic}

\subsection{The Non-Saturating Generator Heuristic}
Early in training, when $G$ generates poor images, $D$ easily rejects them with high confidence ($D(G(z)) \approx 0$). In this regime, the gradient of the standard minimax loss $\log(1 - D(G(z)))$ with respect to $D(G(z))$ is:
\begin{equation}
\frac{d}{dy} \log(1 - y) = -\frac{1}{1 - y} \xrightarrow{y \to 0} -1
\end{equation}
Coupled with the discriminator's sigmoid saturation $\sigma'(a) \to 0$, the gradient transmitted to $G$ vanishes. To overcome this early saturation, Goodfellow et al. proposed the \textit{non-saturating heuristic}, maximizing $\log D(G(z))$ instead of minimizing $\log(1 - D(G(z)))$:
\begin{equation}
\mathcal{L}_G^{NS} = -\mathbb{E}_{z \sim p_z}[\log D(G(z))]
\end{equation}
Its derivative with respect to $y = D(G(z))$ is:
\begin{equation}
\frac{d}{dy} [-\log y] = -\frac{1}{y} \xrightarrow{y \to 0} -\infty
\end{equation}
This provides strong, large gradients when $G$ is performing poorly, accelerating early learning.

\subsection{Pathological Failure Modes}
Despite its theoretical elegance, standard GAN training suffers from acute instability:
\begin{enumerate}
    \item \textbf{Vanishing Gradients:} High-dimensional natural images lie on low-dimensional manifolds embedded in high-dimensional ambient space. When the supports of $p_{data}$ and $p_g$ have disjoint supports or intersect on a set of measure zero, there exists an almost-surely separating hyperplane. Consequently, $D_{JS}(p_{data} \parallel p_g) = \log 2$ everywhere, yielding zero gradient signal to guide $G$.
    \item \textbf{Mode Collapse:} Rather than learning the full multi-modal distribution $p_{data}$, the generator collapses to producing a narrow subset of output variations (or a single mode) that consistently fool $D$. Mathematically, $G$ maximizes reward by placing all probability mass on the single argmax of $D(x)$, ignoring entropy across the latent space $\mathcal{Z}$.
    \item \textbf{Limit Cycles and Non-Convergence:} Simultaneous gradient descent on non-convex minimax objectives is known to cycle or diverge rather than converge, analogous to the vector field differential equations of competitive systems ($\dot{x} = y, \dot{y} = -x$).
\end{enumerate}

\section{Advanced Architectures: Analysis of Seminal Milestones}
To resolve the pathological instabilities of the vanilla GAN, four landmark papers introduced pivotal architectural and mathematical transformations:

\subsection{Milestone 1: DCGAN (Radford et al., 2015)}
Radford, Metz, and Chintala \cite{radford2015unsupervised} bridged convolutional neural networks and generative modeling by introducing Deep Convolutional GANs (DCGAN), establishing concrete architectural design rules:
\begin{itemize}
    \item \textbf{Elimination of Spatial Pooling:} Deterministic spatial pooling (e.g., MaxPool) was replaced entirely by strided convolutions in the Discriminator and fractional-strided (transposed) convolutions in the Generator.
    \item \textbf{Batch Normalization:} Applied throughout both networks to stabilize learning by normalizing activations to zero mean and unit variance, preventing gradient collapse.
    \item \textbf{Activation Topologies:} LeakyReLU (leak slope 0.2) in all discriminator layers to allow gradient flow even for non-activated neurons; ReLU in the generator with Tanh at the output layer.
\end{itemize}

\subsection{Milestone 2: WGAN and WGAN-GP (Arjovsky et al., 2017; Gulrajani et al., 2017)}
Arjovsky et al. \cite{arjovsky2017wasserstein} proved that the fundamental flaw of vanilla GANs stems from the discontinuity of Jensen-Shannon divergence across disjoint manifolds. They replaced JSD with the \textbf{Wasserstein-1 (Earth Mover's) Distance}:
\begin{equation}
W(p_{data}, p_g) = \inf_{\gamma \in \Pi(p_{data}, p_g)} \mathbb{E}_{(x, y) \sim \gamma}[\|x - y\|]
\end{equation}
Using the Kantorovich-Rubinstein duality theorem, the objective becomes:
\begin{equation}
\max_{w \in \mathcal{W}, \|f_w\|_L \le 1} \mathbb{E}_{x \sim p_{data}}[f_w(x)] - \mathbb{E}_{z \sim p_z}[f_w(G(z))]
\end{equation}
where $f_w$ is a critic constrained to the space of 1-Lipschitz continuous functions.

While original WGAN enforced Lipschitz continuity via crude weight clipping ($w \in [-c, c]$), this caused capacity under-utilization and gradient exploding/vanishing. Gulrajani et al. \cite{gulrajani2017improved} solved this via \textbf{Wasserstein GAN with Gradient Penalty (WGAN-GP)}, penalizing critic gradients that deviate from unit norm along straight lines between real and generated points:
\begin{equation}
\mathcal{L}_{WGAN-GP} = \mathbb{E}[\tilde{D}(\tilde{x})] - \mathbb{E}[\tilde{D}(x)] + \lambda \, \mathbb{E}_{\hat{x}}\left[ \left( \|\nabla_{\hat{x}} \tilde{D}(\hat{x})\|_2 - 1 \right)^2 \right]
\end{equation}
where $\hat{x} = \epsilon x + (1 - \epsilon)\tilde{x}$ with $\epsilon \sim \mathcal{U}[0, 1]$, and $\lambda = 10$. This provides continuous, meaningful gradients everywhere, virtually eliminating mode collapse.

\subsection{Milestone 3: Conditional GANs and Image Translation (cGAN / Pix2Pix / CycleGAN)}
Mirza and Osindero \cite{mirza2014conditional} extended GANs to conditional synthesis by feeding side information $y$ (class labels, text embeddings, or conditioning images) to both $G$ and $D$:
\begin{align}
\min_G \max_D V(D, G) = &\;\mathbb{E}_{x, y} [\log D(x \mid y)] \nonumber \\
&+ \mathbb{E}_{z, y} [\log (1 - D(G(z \mid y) \mid y))]
\end{align}
This formulation catalyzed paired image-to-image translation (Pix2Pix \cite{isola2017image}) using U-Net skip architectures with PatchGAN discriminators, and unpaired image-to-image translation (CycleGAN \cite{zhu2017unpaired}) using cycle-consistency loss:
\begin{equation}
\mathcal{L}_{cyc}(G, F) = \mathbb{E}_x[\|F(G(x)) - x\|_1] + \mathbb{E}_y[\|G(F(y)) - y\|_1]
\end{equation}

\subsection{Milestone 4: StyleGAN and StyleGAN2 (Karras et al., 2019/2020)}
Karras et al. \cite{karras2019style, karras2020analyzing} revolutionized high-resolution image synthesis by redesigning the generator architecture based on style transfer principles:
\begin{enumerate}
    \item \textbf{Mapping Network ($f$):} An 8-layer MLP maps latent $z \in \mathcal{Z}$ to intermediate disentangled space $w \in \mathcal{W}$, unconstraining latent representations from the fixed density of $p_z$.
    \item \textbf{Adaptive Instance Normalization (AdaIN):} Vector $w$ is transformed into style scale $y_{s,i}$ and bias $y_{b,i}$ injected into each convolutional block:
    \begin{equation}
    \operatorname{AdaIN}(x_i, y) = y_{s,i} \left( \frac{x_i - \mu(x_i)}{\sigma(x_i)} \right) + y_{b,i}
    \end{equation}
    \item \textbf{Stochastic Variation via Noise Injection:} Explicit Gaussian noise is added per channel, separating coarse global attributes (pose, face shape) from fine stochastic textures (hair strands, pores).
    \item \textbf{StyleGAN2 Weight Demodulation:} StyleGAN2 eliminated water droplet artifacts by baking AdaIN directly into convolution weights via weight demodulation, replacing progressive growing with skip-connected generator architectures.
\end{enumerate}

\begin{table*}[t]
\caption{Comprehensive Comparison of Seminal GAN Architectures}
\label{tab:gan_comparison}
\centering
\small
\begin{tabularx}{\textwidth}{lXXXXX}
\toprule
\textbf{Architecture} & \textbf{Core Mathematical Objective} & \textbf{Divergence / Metric} & \textbf{Lipschitz Constraint} & \textbf{Primary Innovation} & \textbf{Stability / Failure Modes} \\
\midrule
\textbf{Vanilla GAN} \cite{goodfellow2014generative} & $\min_G \max_D \mathbb{E}[\log D] + \mathbb{E}[\log(1-D)]$ & Jensen-Shannon ($D_{JS}$) & None & Minimax game theory foundation & Severe vanishing gradients, mode collapse \\
\textbf{DCGAN} \cite{radford2015unsupervised} & Standard non-saturating $\log D(G(z))$ & Jensen-Shannon ($D_{JS}$) & None (via Batch Normalization) & All-convolutional topology, strided conv & Moderately stable; sensitive to hyper-parameters \\
\textbf{cGAN / Pix2Pix} \cite{mirza2014conditional, isola2017image} & Conditional minimax + $\mathcal{L}_1$ reconstruction & Conditional $D_{JS} + \mathcal{L}_1$ & None & Conditioning vector / PatchGAN discriminator & High conditional fidelity; mode collapse on noise $z$ \\
\textbf{WGAN} \cite{arjovsky2017wasserstein} & $\max_{w} \mathbb{E}[D(x)] - \mathbb{E}[D(G(z))]$ & Earth Mover's ($W_1$) & Weight clipping $[-c, c]$ & Kantorovich-Rubinstein duality & No vanishing gradient; clipping limits capacity \\
\textbf{WGAN-GP} \cite{gulrajani2017improved} & $W_1$ distance + gradient norm penalty & Earth Mover's ($W_1$) & Penalty: $(\|\nabla_{\hat{x}} D\|_2 - 1)^2$ & Soft 1-Lipschitz enforcement everywhere & Highly stable convergence; mode collapse resistant \\
\textbf{StyleGAN / 2} \cite{karras2019style, karras2020analyzing} & Non-saturating GAN + $R_1$ regularization & Disentangled $W_1$ / JSD & $R_1$ gradient penalty & Mapping network $\mathcal{W}$, AdaIN, Demodulation & Exceptional 1024x1024 resolution, fine style control \\
\bottomrule
\end{tabularx}
\end{table*}

\section{Practical Applications and Quantitative Evaluation}

\subsection{Domain Applications}
GAN architectures dominate numerous practical applications:
\begin{itemize}
    \item \textbf{Super-Resolution (SRGAN / ESRGAN):} Combining adversarial loss with perceptual VGG feature loss to hallucinate realistic high-frequency textures in low-resolution imagery \cite{ledig2017photo}.
    \item \textbf{Cross-Domain Synthesis (Pix2Pix / CycleGAN):} Semantic segmentation mask-to-photo synthesis, aerial-to-map cartography, and artistic style transfer.
    \item \textbf{Medical Imaging:} Cross-modality translation (MRI to CT synthesis), reducing radiation exposure and generating synthetic training lesions for rare pathology classifiers.
    \item \textbf{Data Augmentation:} Class-balanced synthetic data generation to overcome severe sample scarcity in industrial defect detection.
\end{itemize}

\subsection{Quantitative Evaluation Metrics}
Evaluating generative models is non-trivial due to the absence of explicit likelihood values. Two dominant metrics govern quantitative benchmarking:
\begin{enumerate}
    \item \textbf{Inception Score (IS):} Measures sharpness and diversity using an Inception-v3 classifier:
    \begin{equation}
    \operatorname{IS}(G) = \exp\left( \mathbb{E}_{x \sim p_g} [D_{KL}(p(y \mid x) \parallel p(y))] \right)
    \end{equation}
    High conditional predictability $p(y \mid x)$ ensures sample clarity, while high marginal entropy $p(y)$ ensures output diversity \cite{salimans2016improved}.
    \item \textbf{Fr\'echet Inception Distance (FID):} Evaluates distribution-level similarity by computing the Wasserstein-2 distance between Gaussian approximations of Inception-v3 pool3 activations for real ($\mu_r, \Sigma_r$) and generated ($\mu_g, \Sigma_g$) sets \cite{heusel2017gans}:
    \begin{equation}
    \operatorname{FID} = \|\mu_r - \mu_g\|_2^2 + \operatorname{Tr}\left(\Sigma_r + \Sigma_g - 2(\Sigma_r \Sigma_g)^{1/2}\right)
    \label{eq:fid}
    \end{equation}
    Lower FID directly correlates with human visual quality and structural diversity.
\end{enumerate}

\section{Conclusion}
Generative Adversarial Networks have fundamentally enriched generative modeling through game-theoretic formulation and implicit density learning. From the initial mathematical derivation of the optimal discriminator and Jensen-Shannon divergence minimization, to the structural rigor of DCGAN, the metric stability of WGAN-GP, and the disentangled generative prowess of StyleGAN, adversarial learning has matured into a powerful mathematical and practical toolkit. Understanding these formal principles is crucial for advancing generative architectures and understanding multi-distribution dynamics in deep learning.

\begin{thebibliography}{99}

\bibitem{goodfellow2014generative}
I.~Goodfellow, J.~Pouget-Abadie, M.~Mirza, B.~Xu, D.~Warde-Farley, S.~Ozair, A.~Courville, and Y.~Bengio, ``Generative adversarial nets,'' in \textit{Advances in Neural Information Processing Systems (NeurIPS)}, vol.~27, pp.~2672--2680, 2014.

\bibitem{radford2015unsupervised}
A.~Radford, L.~Metz, and S.~Chintala, ``Unsupervised representation learning with deep convolutional generative adversarial networks,'' \textit{arXiv preprint arXiv:1511.06434}, 2015.

\bibitem{arjovsky2017wasserstein}
M.~Arjovsky, S.~Chintala, and L.~Bottou, ``Wasserstein generative adversarial networks,'' in \textit{International Conference on Machine Learning (ICML)}, pp.~214--223, 2017.

\bibitem{gulrajani2017improved}
I.~Gulrajani, F.~Ahmed, M.~Arjovsky, V.~Dumoulin, and A.~C.~Courville, ``Improved training of Wasserstein GANs,'' in \textit{Advances in Neural Information Processing Systems (NeurIPS)}, vol.~30, pp.~5767--5777, 2017.

\bibitem{karras2019style}
T.~Karras, S.~Laine, and T.~Aila, ``A style-based generator architecture for generative adversarial networks,'' in \textit{IEEE/CVF Conference on Computer Vision and Pattern Recognition (CVPR)}, pp.~4401--4410, 2019.

\bibitem{karras2020analyzing}
T.~Karras, S.~Laine, M.~Aittala, J.~Hellsten, J.~Lehtinen, and T.~Aila, ``Analyzing and improving the image quality of StyleGAN,'' in \textit{IEEE/CVF Conference on Computer Vision and Pattern Recognition (CVPR)}, pp.~8110--8119, 2020.

\bibitem{mirza2014conditional}
M.~Mirza and S.~Osindero, ``Conditional generative adversarial nets,'' \textit{arXiv preprint arXiv:1411.1784}, 2014.

\bibitem{isola2017image}
P.~Isola, J.-Y.~Zhu, T.~Zhou, and A.~A.~Efros, ``Image-to-image translation with conditional adversarial networks,'' in \textit{IEEE Conference on Computer Vision and Pattern Recognition (CVPR)}, pp.~1125--1134, 2017.

\bibitem{zhu2017unpaired}
J.-Y.~Zhu, T.~Park, P.~Isola, and A.~A.~Efros, ``Unpaired image-to-image translation using cycle-consistent adversarial networks,'' in \textit{IEEE International Conference on Computer Vision (ICCV)}, pp.~2223--2232, 2017.

\bibitem{ledig2017photo}
C.~Ledig, L.~Theis, F.~Husz{\'a}r, J.~Caballero, A.~Cunningham, A.~Acosta, A.~Aitken, A.~Tejani, J.~Totz, Z.~Wang, \textit{et al.}, ``Photo-realistic single image super-resolution using a generative adversarial network,'' in \textit{IEEE Conference on Computer Vision and Pattern Recognition (CVPR)}, pp.~4681--4690, 2017.

\bibitem{salimans2016improved}
T.~Salimans, I.~Goodfellow, W.~Zaremba, V.~Cheung, A.~Radford, and X.~Chen, ``Improved techniques for training GANs,'' in \textit{Advances in Neural Information Processing Systems (NeurIPS)}, vol.~29, pp.~2234--2242, 2016.

\bibitem{heusel2017gans}
M.~Heusel, H.~Ramsauer, T.~Unterthiner, B.~Nessler, and S.~Hochreiter, ``GANs trained by a two time-scale update rule converge to a local Nash equilibrium,'' in \textit{Advances in Neural Information Processing Systems (NeurIPS)}, vol.~30, pp.~6626--6637, 2017.

\end{thebibliography}

\end{document}
